%======================================================================
\chapter{Rods: the elastica and beyond}\label{ch:rd}
%======================================================================
% Outline (23 Nov 2026). Labels of this chapter: *:rd-*  (bibitem keys rd-*).
% Planned content, sources and open questions: ch04_rods/README.md.

\framepart{Outline}
\draftnote{Outline drafted 23 Nov 2026; details, sources and open questions in ch04\_rods/README.md.
Spine (author's slides C05, slide 11): one column, three energies (no rotation, small rotations,
large rotations), three bifurcation diagrams. Energy written $\mathcal{W}$. Prerequisite: a short
section on the bifurcation of a 3D elastic body before this chapter (proposed in Chapter~\ref{ch:nc}).}

\section{From the three-dimensional body to the rod}\label{sec:rd-3d}
\draftnote{Fibre strain $e = u' + \tfrac12 w'^2$ (rotation small but not negligible); Euler--Bernoulli kinematics,
$e = u' + \tfrac12 w'^2 - z w''$; energy
$\mathcal{W} = \tfrac12\int EA(u'+\tfrac12w'^2)^2 + \tfrac12\int EI\, w''^2 - \eta\,\ell$; equilibrium $N'=0$,
$M'' - Nw'' = 0$. Gbox: the energy scale chooses the rod model. Why it works: Ciarlet's asymptotic method
(Berm\'udez--Via\~no, Trabucho--Via\~no; Cimeti\`ere et al.\ for nonlinear rods) and $\Gamma$-convergence
(Mora--M\"uller 2003: elastica; 2004: von K\'arm\'an rod). Cited, not proved.}

\section{Small rotations: Euler loads and design}\label{sec:rd-euler}
\draftnote{Second variation, adjacent equilibrium $EI\,w'''' + P\,w'' = 0$; four supports
($\pi^2EI/L^2$ times $1$, $2.046$, $4$, $1/4$); effective length and Euler hyperbola; flat secondary branches
under a dead load; imperfect column, amplification $1/(P_c/P-1)$, secant formula and Southwell plot;
design limits (Hibbeler Ch.~13).}

\section{Large rotations: the Euler elastica}\label{sec:rd-elastica}
\draftnote{$EI\theta'' + P\sin\theta = 0$ and the pendulum analogy (Audoly--Pomeau Ch.~3, Triantafyllidis BA-1);
closed form with $k=\sin(\theta_0/2)$, $P/P_c = (2K(k)/\pi)^2$; table of shapes; Koiter by Box~2.2,
$P/P_c = 1 + \theta_0^2/8$: stable-symmetric but weak; higher modes unstable.}

\section{Beams on an elastic foundation: wavelength selection}\label{sec:rd-foundation}
\draftnote{$EI\,w'''' + P\,w'' + k\,w = 0$; $P_c = 2\sqrt{kEI}$, wavelength $2\pi(EI/k)^{1/4}$ independent of
the length; finite length and mode jumps; rails, pipelines, film on substrate. With the elastica the periodic
branch is unstable, $P/P_c \approx 1 - A^2/8$ (to confirm); localisation cited.}

\section{Twist and natural curvature}\label{sec:rd-twist}
\draftnote{Short, or moved to Chapter~8 (open): Kirchhoff rod in 3D, Greenhill's shaft, localised helix,
naturally curved rods.}

\framepart{Summary}
\begin{gbox*}
\begin{itemize}
  \item \textbf{To be written.} \draftnote{Rotations decide the post-buckling; Euler loads and effective length;
  the elastica is stable but weak; a foundation picks the wavelength and is unstable; trust a code after the elastica.}
\end{itemize}
\end{gbox*}

\framepart{Exercises}\label{sec:rd-exo}
\framesub{Short questions}
\draftnote{Effective lengths, flat branch, Southwell, pendulum analogy, rail wavelength, $h^2$ against $h^4$,
yield or buckle (SI units).}
\framesub{Problems}
\draftnote{Triantafyllidis H3, H4; clamped--free elastica; finite beam on a foundation; secant-formula design (Hibbeler).}
\framesub{Computer exercises}
\draftnote{Elastica by shooting and continuation (Python); elliptic solution (Mathematica); FEniCS and Cast3M
against the closed form.}

\begin{chapterbib}{99}
\bibitem{rd-audoly2010} B.~Audoly and Y.~Pomeau, \emph{Elasticity and Geometry: From Hair Curls to the Non-linear Response of Shells}, Oxford University Press, 2010.
\bibitem{rd-timoshenko1961} S.~P.~Timoshenko and J.~M.~Gere, \emph{Theory of Elastic Stability}, 2nd ed., McGraw-Hill, New York, 1961.
\bibitem{rd-bergou2008} M.~Bergou, M.~Wardetzky, S.~Robinson, B.~Audoly and E.~Grinspun, Discrete elastic rods, \emph{ACM Transactions on Graphics} 27 (2008) 63.
\bibitem{rd-reis2015} P.~M.~Reis, A perspective on the revival of structural (in)stability with novel opportunities for function: from buckliphobia to buckliphilia, \emph{Journal of Applied Mechanics} 82 (2015) 111001.
\bibitem{rd-triantafyllidis2011} N.~Triantafyllidis, \emph{Stability of Solids: from Structures to Materials}, lecture notes MEC563, \'Ecole polytechnique, Palaiseau, 2011.
\end{chapterbib}
